\section{Stereographic Block Encoding}
\label{sec:block-encoding}

The single-qubit theory of
Section~\ref{sec:review} transforms a scalar signal $r$ into
an unbounded rational function.  To apply this machinery to a
multi-qubit Hamiltonian
$H = \sum_j \lambda_j\ket{j}\bra{j}$, one needs a way to feed
each eigenvalue $\lambda_j$ into the QSP circuit
simultaneously.  This section defines the stereographic block
encoding, compares it with the standard framework, and analyzes
its structural properties.

\subsection{Standard Block Encoding}
\label{subsec:standard-block}

To fix notation and set up the comparison, we begin with the
standard framework.

\begin{definition}[Standard Block Encoding~\cite{GSLW2019}]
	\label{def:standard-block}
	An $(\alpha, m, \epsilon)$-block encoding of a Hermitian
	operator $H$ acting on $n$ qubits is a unitary $U$ acting
	on $n + m$ qubits such that
	\begin{equation}
		\|H - \alpha(\bra{0}^{\otimes m} \otimes I)\, U\,
		(\ket{0}^{\otimes m} \otimes I)\| \leq \epsilon,
		\label{eq:standard-block}
	\end{equation}
	where $\alpha \geq \|H\|$.
\end{definition}

To apply QSVT and compute $f(H)$, one constructs a degree-$d$
polynomial $P$ satisfying $|P(\lambda/\alpha)| \leq 1$ for all
eigenvalues $\lambda$, then finds QSP phases for $P$.  This
approach incurs two costs: the subnormalization factor
$\alpha \geq \|H\|$ increases circuit depth by a factor
of~$\alpha$, and the polynomial $P$ must approximate
$f(\alpha\,\cdot\,)$ on $[-1,1]$, which may require high
degree if $f$ has sharp features or singularities.

\subsection{Stereographic Block Encoding}
\label{subsec:stereo-block}

The stereographic alternative removes both costs by encoding
each eigenvalue directly, without normalization.

\begin{definition}[Stereographic Block Encoding]
	\label{def:stereo-block}
	Let $H$ be a Hermitian operator with spectral decomposition
	$H = \sum_j \lambda_j \ket{j}\bra{j}$, where
	$\lambda_j \geq 0$.  A stereographic block encoding of~$H$
	is a unitary $U$ acting on $n + 1$ qubits such that
	\begin{equation}
		U = \sum_j U_{\lambda_j} \otimes \ket{j}\bra{j},
		\label{eq:stereo-block}
	\end{equation}
	where $U_{\lambda_j}$ is the single-qubit encoding
	unitary~\eqref{eq:encoding-unitary} with $r = \lambda_j$
	and $\theta = 0$.
\end{definition}

Explicitly, in the eigenspace of eigenvalue $\lambda_j$:
\begin{equation}
	U_{\lambda_j} = \frac{1}{\sqrt{1+\lambda_j^2}}
	\begin{pmatrix} \lambda_j & 1 \\ 1 & -\lambda_j \end{pmatrix},
	\label{eq:encoding-eigenspace}
\end{equation}
which maps $\ket{0}\ket{j} \mapsto \ket{\lambda_j}\ket{j}$
using a single ancilla qubit.

Several structural properties distinguish the stereographic
construction from the standard one.

\begin{proposition}[Properties of Stereographic Block Encoding]
	\label{prop:stereo-properties}
	Let $U$ be a stereographic block encoding of~$H$ as in
	Definition~\ref{def:stereo-block}.  Then:
	\begin{enumerate}
		\item \textbf{Exactness.}  The encoding is exact:
		      there is no approximation error $\epsilon$.
		      Each eigenvalue $\lambda_j$ is encoded into
		      the Bloch sphere coordinates without
		      truncation.
		\item \textbf{Single ancilla.}  The encoding uses
		      exactly one ancilla qubit, regardless of
		      $\|H\|$ or the number of system qubits~$n$.
		\item \textbf{No normalization.}  The encoding does
		      not require knowledge of $\|H\|$, and no
		      normalization factor enters the construction.
		\item \textbf{Eigenbasis dependence.}
		      Constructing~\eqref{eq:stereo-block} requires
		      the ability to apply $U_{\lambda_j}$
		      conditioned on the eigenspace $\ket{j}$
		      of~$H$.
	\end{enumerate}
\end{proposition}

\begin{proof}
	Parts (1)--(3) follow directly from the
	definition: $U_{\lambda_j}$ is the exact encoding
	unitary for $r = \lambda_j$, acts on a single ancilla
	qubit, and depends only on~$\lambda_j$ (not on any
	global property of~$H$).  Part~(4) is immediate from
	the structure of~\eqref{eq:stereo-block}: the sum over
	eigenspace projectors $\ket{j}\bra{j}$ requires the
	ability to distinguish eigenspaces.
\end{proof}

\subsection{The Eigenbasis Requirement}
\label{subsec:eigenbasis}

The eigenbasis requirement (Proposition~\ref{prop:stereo-properties},
part~4) is the principal limitation of stereographic block encoding
and deserves careful analysis.  In the standard QSVT framework,
$H/\alpha$ is embedded into a unitary without knowledge of the
eigenstates, and the QSP polynomial acts simultaneously in all
eigenspaces precisely because the block encoding does not
distinguish them.  The eigenbasis requirement restricts
stereographic block encoding to \emph{structured Hamiltonians}
whose eigenbasis is known or efficiently computable.

We identify three classes of increasing generality:

\textit{Class~1: Computational-basis diagonal operators.}
Hamiltonians that are diagonal in the computational basis require
no change of basis: the encoding rotation is conditioned directly
on the computational-basis state.  This includes classical Ising
models, diagonal parts of molecular Hamiltonians, and any operator
of the form $H = \sum_j \lambda_j \ket{j}\bra{j}$ with $\ket{j}$
being computational-basis states.

\textit{Class~2: Known diagonalization.}
Hamiltonians $H = V^\dagger D V$ where the change-of-basis
unitary~$V$ is efficiently implementable as a quantum circuit.
This includes integrable spin chains (solved by Jordan--Wigner or
Bethe ansatz), free-fermion models (Bogoliubov transformation),
and Hamiltonians with analytically known eigenbases (e.g., the
Heisenberg model on small systems).

\textit{Class~3: QPE-assisted encoding.}  For Hamiltonians where
the eigenbasis is not known a priori, quantum phase estimation
(QPE) can in principle extract eigenvalues into a register,
enabling conditioned rotations.  However, this reintroduces the
resources (ancilla qubits, circuit depth) that stereographic
encoding aims to avoid, and the advantage over standard QSVT
becomes unclear.  We discuss this connection further in
Section~\ref{sec:discussion}.

\begin{remark}[Positivity assumption]
	\label{rem:positivity}
	Definition~\ref{def:stereo-block} assumes
	$\lambda_j \geq 0$.  For Hamiltonians with negative
	eigenvalues, a spectral shift
	$H' = H + |\lambda_{\min}|I \geq 0$ restores positivity.
	The shift changes the target function:
	$f(\lambda) \to f(\lambda - |\lambda_{\min}|)$, which must
	be accounted for in the polynomial design.  The shift
	$|\lambda_{\min}|$ is a single scalar and does not require
	knowledge of the full eigenbasis.
\end{remark}

\subsection{Comparison with Standard Block Encoding}
\label{subsec:comparison}

Table~\ref{tab:block-comparison} summarizes the structural
differences.

\begin{table}[t]
	\caption{Comparison of standard and stereographic block
		encoding for a Hermitian operator $H$ with
		$\|H\| = \alpha$.}
	\label{tab:block-comparison}
	\begin{ruledtabular}
		\begin{tabular}{lll}
			                 & Standard              & Stereographic    \\
			\hline
			Encoding type    & Approximate           & Exact            \\
			Ancilla qubits   & $m \geq 1$            & $1$              \\
			Normalization    & $\alpha \geq \|H\|$   & Not required     \\
			Eigenbasis       & Not required          & Required         \\
			Depth per query  & $\mathcal{O}(\alpha)$ & $\mathcal{O}(1)$ \\
			Applicable class & All Hermitian         & Structured       \\
		\end{tabular}
	\end{ruledtabular}
\end{table}
